% =====================================================================
% Abbildung: TF-Baum und physische Bedeutung der Koordinatenframes
% Passend zu Kapitel 10, Abschnitt "Koordinatensysteme"
%
% Einbindung im Buch, z. B. in chapters/10_schnittstellen.tex:
%
%   \begin{figure}[H]
%     \centering
%     % =====================================================================
% Abbildung: TF-Baum und physische Bedeutung der Koordinatenframes
% Passend zu Kapitel 10, Abschnitt "Koordinatensysteme"
%
% Einbindung im Buch, z. B. in chapters/10_schnittstellen.tex:
%
%   \begin{figure}[H]
%     \centering
%     % =====================================================================
% Abbildung: TF-Baum und physische Bedeutung der Koordinatenframes
% Passend zu Kapitel 10, Abschnitt "Koordinatensysteme"
%
% Einbindung im Buch, z. B. in chapters/10_schnittstellen.tex:
%
%   \begin{figure}[H]
%     \centering
%     % =====================================================================
% Abbildung: TF-Baum und physische Bedeutung der Koordinatenframes
% Passend zu Kapitel 10, Abschnitt "Koordinatensysteme"
%
% Einbindung im Buch, z. B. in chapters/10_schnittstellen.tex:
%
%   \begin{figure}[H]
%     \centering
%     \input{figures/fig_kap10_tf_baum}
%     \caption{Links die Hierarchie der TF-Frames vom kartenfesten
%       \texttt{map} über \texttt{odom} und \texttt{base\_link} bis zu
%       \texttt{camera\_link}; rechts dieselben Frames als physische
%       Koordinatenursprünge im Gebäudegrundriss und am Roboter.}
%     \label{fig:kap10-tf-baum}
%   \end{figure}
%
% Benoetigte Pakete (im Buch bereits in main.tex geladen): tikz
% Benoetigte TikZ-Bibliotheken: positioning, arrows.meta, shadows, calc
% =====================================================================

\input{figures/fig_kap10_tf_baum.tikz}

%     \caption{Links die Hierarchie der TF-Frames vom kartenfesten
%       \texttt{map} über \texttt{odom} und \texttt{base\_link} bis zu
%       \texttt{camera\_link}; rechts dieselben Frames als physische
%       Koordinatenursprünge im Gebäudegrundriss und am Roboter.}
%     \label{fig:kap10-tf-baum}
%   \end{figure}
%
% Benoetigte Pakete (im Buch bereits in main.tex geladen): tikz
% Benoetigte TikZ-Bibliotheken: positioning, arrows.meta, shadows, calc
% =====================================================================

\begin{tikzpicture}[
    every node/.style={font=\small},
    tf/.style={
        rectangle, rounded corners=2pt, draw=violet!55!black, thick,
        minimum width=24mm, minimum height=9mm,
        align=center, fill=violet!12, drop shadow
    },
    flow/.style={-{Stealth[length=2mm]}, thick, violet!55!black},
    lbl/.style={font=\tiny, align=left, text width=21mm},
    room/.style={draw=black!50, thick},
    axis/.style={-{Stealth[length=1.4mm]}, thick},
    origin/.style={circle, fill=black, minimum size=1.2mm, inner sep=0pt}
]

% ===== Linke Seite: TF-Baum (Hierarchie) =====
\node[font=\small\bfseries, align=left, text width=48mm] (title1) at (1.3,10.6) {TF-Baum\\(Hierarchie der Frames)};

\node[tf] (map)  at (1.3,9)  {\texttt{map}};
\node[tf] (odom) at (1.3,6.2) {\texttt{odom}};
\node[tf] (base) at (1.3,3.4) {\texttt{base\_link}};
\node[tf] (cam)  at (1.3,0.6) {\texttt{camera\_link}};

\draw[flow] (map.south)  -- (odom.north);
\node[lbl] at (4.1,7.6) {SLAM-/AMCL-Korrektur (springt gelegentlich)};

\draw[flow] (odom.south) -- (base.north);
\node[lbl] at (4.1,4.8) {Odometrie aus Rad\-encodern/IMU (driftet langsam)};

\draw[flow] (base.south) -- (cam.north);
\node[lbl] at (4.1,2.0) {statische Transformation (fest montiert)};

% ===== Trennlinie =====
\draw[black!30, thick] (6.6,11.1) -- (6.6,-1.3);

% ===== Rechte Seite: physische Bedeutung (Draufsicht) =====
\node[font=\small\bfseries, align=left, text width=55mm] (title2) at (10.5,10.6) {Physische Bedeutung\\(Draufsicht)};

\draw[room] (7,-0.8) rectangle (14.5,9.3);
\node[font=\tiny, black!50] at (8.4,8.9) {Gebäudegrundriss};

% map-Ursprung: fest, z. B. Kartenursprung/Startpunkt
\node[origin] (mo) at (7.9,0.4) {};
\draw[axis, red!70!black]   (mo) -- ++(0.8,0)  node[right, font=\tiny]{x};
\draw[axis, green!50!black] (mo) -- ++(0,0.8)  node[above, font=\tiny]{y};
\node[font=\tiny, align=left, text width=20mm] at (7.9,-1.6) {\texttt{map} (fest, Ursprung der Karte)};

% odom-Ursprung: nahe map, aber mit Drift (gestrichelter Kreis)
\node[origin, fill=violet!55!black] (oo) at (8.5,1.0) {};
\draw[dashed, violet!45] (8.5,1.0) circle (3.5mm);
\node[font=\tiny, violet!55!black, align=left, text width=20mm] at (11.3,0.3) {\texttt{odom} (driftet mit der Zeit gegenüber \texttt{map})};

% Roboter = base_link
\node[tf, minimum width=18mm, minimum height=13mm] (robot) at (11.3,5.2) {};
\node[font=\tiny\bfseries, black!70] at (11.3,6.9) {Roboter};
\node[origin] at (robot.center) {};
\draw[axis, red!70!black]   (robot.center) -- ++(0.55,0);
\draw[axis, green!50!black] (robot.center) -- ++(0,0.55);
\node[font=\tiny, align=center, text width=20mm] at (11.3,3.3) {\texttt{base\_link} (Robotermitte)};

% Kamera = camera_link, vorne am Roboter versetzt
\node[fill=orange!15, draw=orange!60!black, thick, rounded corners=1pt,
      minimum width=7mm, minimum height=5mm] (camnode) at (13.6,6.6) {};
\node[origin] at (camnode.center) {};
\draw[black!40, thick] (robot.north east) -- (camnode.south west);
\node[font=\tiny, align=left, text width=19mm] at (13.0,8.3) {\texttt{camera\_link} (Halterung, ca. 15\,cm vor Mitte)};

\end{tikzpicture}


%     \caption{Links die Hierarchie der TF-Frames vom kartenfesten
%       \texttt{map} über \texttt{odom} und \texttt{base\_link} bis zu
%       \texttt{camera\_link}; rechts dieselben Frames als physische
%       Koordinatenursprünge im Gebäudegrundriss und am Roboter.}
%     \label{fig:kap10-tf-baum}
%   \end{figure}
%
% Benoetigte Pakete (im Buch bereits in main.tex geladen): tikz
% Benoetigte TikZ-Bibliotheken: positioning, arrows.meta, shadows, calc
% =====================================================================

\begin{tikzpicture}[
    every node/.style={font=\small},
    tf/.style={
        rectangle, rounded corners=2pt, draw=violet!55!black, thick,
        minimum width=24mm, minimum height=9mm,
        align=center, fill=violet!12, drop shadow
    },
    flow/.style={-{Stealth[length=2mm]}, thick, violet!55!black},
    lbl/.style={font=\tiny, align=left, text width=21mm},
    room/.style={draw=black!50, thick},
    axis/.style={-{Stealth[length=1.4mm]}, thick},
    origin/.style={circle, fill=black, minimum size=1.2mm, inner sep=0pt}
]

% ===== Linke Seite: TF-Baum (Hierarchie) =====
\node[font=\small\bfseries, align=left, text width=48mm] (title1) at (1.3,10.6) {TF-Baum\\(Hierarchie der Frames)};

\node[tf] (map)  at (1.3,9)  {\texttt{map}};
\node[tf] (odom) at (1.3,6.2) {\texttt{odom}};
\node[tf] (base) at (1.3,3.4) {\texttt{base\_link}};
\node[tf] (cam)  at (1.3,0.6) {\texttt{camera\_link}};

\draw[flow] (map.south)  -- (odom.north);
\node[lbl] at (4.1,7.6) {SLAM-/AMCL-Korrektur (springt gelegentlich)};

\draw[flow] (odom.south) -- (base.north);
\node[lbl] at (4.1,4.8) {Odometrie aus Rad\-encodern/IMU (driftet langsam)};

\draw[flow] (base.south) -- (cam.north);
\node[lbl] at (4.1,2.0) {statische Transformation (fest montiert)};

% ===== Trennlinie =====
\draw[black!30, thick] (6.6,11.1) -- (6.6,-1.3);

% ===== Rechte Seite: physische Bedeutung (Draufsicht) =====
\node[font=\small\bfseries, align=left, text width=55mm] (title2) at (10.5,10.6) {Physische Bedeutung\\(Draufsicht)};

\draw[room] (7,-0.8) rectangle (14.5,9.3);
\node[font=\tiny, black!50] at (8.4,8.9) {Gebäudegrundriss};

% map-Ursprung: fest, z. B. Kartenursprung/Startpunkt
\node[origin] (mo) at (7.9,0.4) {};
\draw[axis, red!70!black]   (mo) -- ++(0.8,0)  node[right, font=\tiny]{x};
\draw[axis, green!50!black] (mo) -- ++(0,0.8)  node[above, font=\tiny]{y};
\node[font=\tiny, align=left, text width=20mm] at (7.9,-1.6) {\texttt{map} (fest, Ursprung der Karte)};

% odom-Ursprung: nahe map, aber mit Drift (gestrichelter Kreis)
\node[origin, fill=violet!55!black] (oo) at (8.5,1.0) {};
\draw[dashed, violet!45] (8.5,1.0) circle (3.5mm);
\node[font=\tiny, violet!55!black, align=left, text width=20mm] at (11.3,0.3) {\texttt{odom} (driftet mit der Zeit gegenüber \texttt{map})};

% Roboter = base_link
\node[tf, minimum width=18mm, minimum height=13mm] (robot) at (11.3,5.2) {};
\node[font=\tiny\bfseries, black!70] at (11.3,6.9) {Roboter};
\node[origin] at (robot.center) {};
\draw[axis, red!70!black]   (robot.center) -- ++(0.55,0);
\draw[axis, green!50!black] (robot.center) -- ++(0,0.55);
\node[font=\tiny, align=center, text width=20mm] at (11.3,3.3) {\texttt{base\_link} (Robotermitte)};

% Kamera = camera_link, vorne am Roboter versetzt
\node[fill=orange!15, draw=orange!60!black, thick, rounded corners=1pt,
      minimum width=7mm, minimum height=5mm] (camnode) at (13.6,6.6) {};
\node[origin] at (camnode.center) {};
\draw[black!40, thick] (robot.north east) -- (camnode.south west);
\node[font=\tiny, align=left, text width=19mm] at (13.0,8.3) {\texttt{camera\_link} (Halterung, ca. 15\,cm vor Mitte)};

\end{tikzpicture}


%     \caption{Links die Hierarchie der TF-Frames vom kartenfesten
%       \texttt{map} über \texttt{odom} und \texttt{base\_link} bis zu
%       \texttt{camera\_link}; rechts dieselben Frames als physische
%       Koordinatenursprünge im Gebäudegrundriss und am Roboter.}
%     \label{fig:kap10-tf-baum}
%   \end{figure}
%
% Benoetigte Pakete (im Buch bereits in main.tex geladen): tikz
% Benoetigte TikZ-Bibliotheken: positioning, arrows.meta, shadows, calc
% =====================================================================

\begin{tikzpicture}[
    every node/.style={font=\small},
    tf/.style={
        rectangle, rounded corners=2pt, draw=violet!55!black, thick,
        minimum width=24mm, minimum height=9mm,
        align=center, fill=violet!12, drop shadow
    },
    flow/.style={-{Stealth[length=2mm]}, thick, violet!55!black},
    lbl/.style={font=\tiny, align=left, text width=21mm},
    room/.style={draw=black!50, thick},
    axis/.style={-{Stealth[length=1.4mm]}, thick},
    origin/.style={circle, fill=black, minimum size=1.2mm, inner sep=0pt}
]

% ===== Linke Seite: TF-Baum (Hierarchie) =====
\node[font=\small\bfseries, align=left, text width=48mm] (title1) at (1.3,10.6) {TF-Baum\\(Hierarchie der Frames)};

\node[tf] (map)  at (1.3,9)  {\texttt{map}};
\node[tf] (odom) at (1.3,6.2) {\texttt{odom}};
\node[tf] (base) at (1.3,3.4) {\texttt{base\_link}};
\node[tf] (cam)  at (1.3,0.6) {\texttt{camera\_link}};

\draw[flow] (map.south)  -- (odom.north);
\node[lbl] at (4.1,7.6) {SLAM-/AMCL-Korrektur (springt gelegentlich)};

\draw[flow] (odom.south) -- (base.north);
\node[lbl] at (4.1,4.8) {Odometrie aus Rad\-encodern/IMU (driftet langsam)};

\draw[flow] (base.south) -- (cam.north);
\node[lbl] at (4.1,2.0) {statische Transformation (fest montiert)};

% ===== Trennlinie =====
\draw[black!30, thick] (6.6,11.1) -- (6.6,-1.3);

% ===== Rechte Seite: physische Bedeutung (Draufsicht) =====
\node[font=\small\bfseries, align=left, text width=55mm] (title2) at (10.5,10.6) {Physische Bedeutung\\(Draufsicht)};

\draw[room] (7,-0.8) rectangle (14.5,9.3);
\node[font=\tiny, black!50] at (8.4,8.9) {Gebäudegrundriss};

% map-Ursprung: fest, z. B. Kartenursprung/Startpunkt
\node[origin] (mo) at (7.9,0.4) {};
\draw[axis, red!70!black]   (mo) -- ++(0.8,0)  node[right, font=\tiny]{x};
\draw[axis, green!50!black] (mo) -- ++(0,0.8)  node[above, font=\tiny]{y};
\node[font=\tiny, align=left, text width=20mm] at (7.9,-1.6) {\texttt{map} (fest, Ursprung der Karte)};

% odom-Ursprung: nahe map, aber mit Drift (gestrichelter Kreis)
\node[origin, fill=violet!55!black] (oo) at (8.5,1.0) {};
\draw[dashed, violet!45] (8.5,1.0) circle (3.5mm);
\node[font=\tiny, violet!55!black, align=left, text width=20mm] at (11.3,0.3) {\texttt{odom} (driftet mit der Zeit gegenüber \texttt{map})};

% Roboter = base_link
\node[tf, minimum width=18mm, minimum height=13mm] (robot) at (11.3,5.2) {};
\node[font=\tiny\bfseries, black!70] at (11.3,6.9) {Roboter};
\node[origin] at (robot.center) {};
\draw[axis, red!70!black]   (robot.center) -- ++(0.55,0);
\draw[axis, green!50!black] (robot.center) -- ++(0,0.55);
\node[font=\tiny, align=center, text width=20mm] at (11.3,3.3) {\texttt{base\_link} (Robotermitte)};

% Kamera = camera_link, vorne am Roboter versetzt
\node[fill=orange!15, draw=orange!60!black, thick, rounded corners=1pt,
      minimum width=7mm, minimum height=5mm] (camnode) at (13.6,6.6) {};
\node[origin] at (camnode.center) {};
\draw[black!40, thick] (robot.north east) -- (camnode.south west);
\node[font=\tiny, align=left, text width=19mm] at (13.0,8.3) {\texttt{camera\_link} (Halterung, ca. 15\,cm vor Mitte)};

\end{tikzpicture}

